% !TeX spellcheck = hu_HU
% !TeX encoding = UTF-8
% !TeX program = xelatex
%----------------------------------------------------------------------------
\section{Argumentum-kötött függvények}\label{sec:ftr:argfun}
%----------------------------------------------------------------------------

A C4 nyelv által biztosított \gls{DSL} leírási lehetőség a megfelelő függvények definiálásával történik.
Ez az általános megoldás majdnem teljesen szabadon formálhatóvá teszik a nyelvet a \gls{DSL} írója számára, aki a saját feladatára személyre tudja szabni, hogy a \gls{DSL} számára megfelelő legyen a szintaktika.

Egy kiemelkedően jól működő struktúra a beágyazott \gls{DSL}-ek megvalósítására olyan függvények bevezetése, amelyek megfelelő módon hívják egymást, és a szintaktikai elemek megadó külső függvények csak ezeket állítják elő valamilyen closure formájában.
Ezeknek a függvényeknek pedig van egy kezdő függvénye, amelyik bevezeti az éppen használt \gls{DSL}-t.
Ez a módszer látható az előző fejezetben adott példák C4 kódjában, alább pedig kiemelve az állapotgép leíró példából a {\tt state} függvény megvalósítása.

\begin{figure}[htb]
	\begin{subfigure}{.45\linewidth}
		\begin{lstlisting}[numbers=left,label=lst:fsm-state-fn]
let state/2 { |name exec|
    { |sm last str err fin|
        let sm2 add_state sm name str err fin
        &(exec)/5 sm2 name
                  nil_block nil_block nil_block
    }
}\end{lstlisting}
		\caption{Az állapotgép \gls{DSL} {\tt state} függvénye}
	\end{subfigure}
	\hfill
	\begin{subfigure}{.45\linewidth}
		\begin{lstlisting}[numbers=left,label=lst:fsm-state-boundfn]
let [fsm/2 state/2] { |name exec|
    { |sm last str err fin|
        let sm2 add_state sm name str err fin
        &(exec)/5 sm2 name
                  nil_block nil_block nil_block
    }
}\end{lstlisting}
		\caption{Az argumentum-kötöt {\tt state} függvény}
	\end{subfigure}
	\caption{A jelenlegi és eltervezett megvalósítás argumentum-kötött függvénnyel}
\end{figure}	

Ennek a struktúrának előnye, hogy egy-egy \gls{DSL} használata teljesen egy függvényhívásba limitálódik, hiszen maga a \gls{DSL}-t megalkotó függvények is csak ezen a függvényen belül fognak az elvárt módon viselkedni.

Ezt a gondolatot megerősíteni kívánó egyik fejlesztési terv olyan függvények bevezetése, amelyek csak más függvények argumentumaként állhatnak.
Ezzel megoldódik az esetleges \gls{DSL}-ek közötti névütközés, hiszen egy-egy kulcsszó csak a megfelelő \gls{DSL}-en belül kerül feloldásra, így minden \gls{DSL} a saját verzióját fogja megkapni.
Ezeket a függvényeket nevezem argumentum-kötött (argument-bound) függvényeknek, mert meg van kötve, hogy melyik másik függvény argumentumában kerülhetnek meghívásra.

Fontos megjegyezni, hogy a tartalmazás a jelenlegi koncepció szerint rekurzívan van értelmezve, miszerint a hívott függvények sorában valahol egyszer meg kell jelennie a várt függvénynek.
Ennek speciális esetét, ahol szigorúan a hívott függvény argumentumaként kell állnia is érdemes mérlegelni. Megéri-e az a szigorúság, amit bevezet egy alapból nagyon megengedő nyelvbe?















